\section{Pillar 1 --- Resource Contention and Performance Collapse}
\textit{(Theoretical basis for \sq{1})}

Resource contention arises when multiple concurrent tasks compete for finite shared
resources --- most commonly \textsc{cpu} cycles and memory --- on a single host. Unlike
bandwidth-limited degradation, resource contention causes non-linear throughput collapse as
load approaches system capacity \citep{nanda1991}.

\citet{nanda1991} demonstrate that as resource utilisation approaches 100\%, response time
grows super-linearly. For a Docker executor on a single VM, this means that adding one additional
concurrent job container beyond the system's stable operating region does not produce a
proportional decrease in throughput --- it produces a disproportionately large degradation,
which can cascade into complete system failure.

In the context of workflow orchestration, this manifests as individual job containers failing
due to \textsc{cpu} starvation or memory exhaustion through the Docker daemon, which can
cause the orchestration process itself to become unresponsive, delaying or failing other
running containers. This is the failure mode that \sq{1} empirically measures and quantifies
for the specific workload profiles defined in this study.

\textbf{Research connection:} \sq{1} identifies the empirical concurrency threshold at
which non-linear collapse occurs for \dagster{} pipeline workloads on a single \vm{}.
